\documentclass[10pt,a4paper]{amsart}
\usepackage{amsmath,amssymb,amsthm,graphicx}
% For commutative diagrams we use Xy-pic
\usepackage[all]{xy}

% We use 2cell for 2-commutative diagrams.
\xyoption{2cell}
% Theorem environments.
%
\theoremstyle{plain}
\newtheorem{theorem}[subsection]{Theorem}
\newtheorem{proposition}[subsection]{Proposition}
\newtheorem{lemma}[subsection]{Lemma}

\theoremstyle{definition}
\newtheorem{definition}[subsection]{Definition}
\newtheorem{example}[subsection]{Example}
\newtheorem{exercise}[subsection]{Exercise}
\newtheorem{situation}[subsection]{Situation}

\theoremstyle{remark}
\newtheorem{remark}[subsection]{Remark}
\newtheorem{remarks}[subsection]{Remarks}

\numberwithin{equation}{subsection}

% Macros
%
\def\lim{\mathop{\mathrm{lim}}\nolimits}
\def\colim{\mathop{\mathrm{colim}}\nolimits}
\def\Spec{\mathop{\mathrm{Spec}}}
\def\Hom{\mathop{\mathrm{Hom}}\nolimits}
\def\Ext{\mathop{\mathrm{Ext}}\nolimits}
\def\SheafHom{\mathop{\mathcal{H}\!\mathit{om}}\nolimits}
\def\SheafExt{\mathop{\mathcal{E}\!\mathit{xt}}\nolimits}
\def\Sch{\mathit{Sch}}
\def\Mor{\mathop{\mathrm{Mor}}\nolimits}
\def\Ob{\mathop{\mathrm{Ob}}\nolimits}
\def\Sh{\mathop{\mathit{Sh}}\nolimits}
\def\NL{\mathop{N\!L}\nolimits}
\def\CH{\mathop{\mathrm{CH}}\nolimits}
\def\proetale{{pro\text{-}\acute{e}tale}}
\def\etale{{\acute{e}tale}}
\def\QCoh{\mathit{QCoh}}
\def\Ker{\mathop{\mathrm{Ker}}}
\def\Im{\mathop{\mathrm{Im}}}
\def\Coker{\mathop{\mathrm{Coker}}}
\def\Coim{\mathop{\mathrm{Coim}}}

% Boxtimes
%
\DeclareMathSymbol{\boxtimes}{\mathbin}{AMSa}{"02}

%
% Macros for moduli stacks/spaces
%
\def\QCohstack{\mathcal{QC}\!\mathit{oh}}
\def\Cohstack{\mathcal{C}\!\mathit{oh}}
\def\Spacesstack{\mathcal{S}\!\mathit{paces}}
\def\Quotfunctor{\mathrm{Quot}}
\def\Hilbfunctor{\mathrm{Hilb}}
\def\Curvesstack{\mathcal{C}\!\mathit{urves}}
\def\Polarizedstack{\mathcal{P}\!\mathit{olarized}}
\def\Complexesstack{\mathcal{C}\!\mathit{omplexes}}
% \Pic is the operator that assigns to X its picard group, usage \Pic(X)
% \Picardstack_{X/B} denotes the Picard stack of X over B
% \Picardfunctor_{X/B} denotes the Picard functor of X over B
\def\Pic{\mathop{\mathrm{Pic}}\nolimits}
\def\Picardstack{\mathcal{P}\!\mathit{ic}}
\def\Picardfunctor{\mathrm{Pic}}
\def\Deformationcategory{\mathcal{D}\!\mathit{ef}}
\usepackage{cleveref}
\crefname{originaldefinition}{original definition}{original definitions}
\Crefname{originaldefinition}{Original definition}{Original definitions}
\makeatletter
\newlabel{sites-modules-definition-g-shriek}{{04BF}{source}{}{sourcepointer-1}{}}
\newlabel{sites-modules-definition-g-shriek@cref}{{[originaldefinition][0][]04BF}{[1][0][]source}{}{}{}}
\makeatother
\title{Comparing structure sheaves for nested full subcategories of the big site}
\author{The Stacks Project Authors}
\date{stacks-project:etale-cohomology.tex:local-snapshot-20261011; GNU FDL 1.2 or later}
\begin{document}
\maketitle
\noindent\textit{Editorial note.} All delivered mathematical text is the author's own, copied from the pinned Stacks Project chapter file etale-cohomology.tex; no statement or proof step was rewritten, reordered or removed.\par\medskip
\noindent\textit{Editorial note.} The proof is the author's complete argument for both parts: surjectivity and injectivity of g\_!O\_1 -> O\_2 are each written out, including the hyperplane trick used for injectivity.\par\medskip
\noindent\textit{Editorial note.} The label lemma-compare-structure-sheaves occurs in no other staged Stacks chapter, so no cross-chapter restatement had to be weighed for this card.\par\medskip
\noindent\textit{Editorial note.} The single labelled reference occurring inside the counted statement points to a chapter of the Project not staged in this card's source tree; it is carried in assembly.external\_references as a non-load-bearing navigation pointer recorded at the invoking line.\par\medskip
\noindent\textit{Editorial note.} The preamble is the chapter's own native preamble excerpt already used for this source: the Xy-pic setup (native preamble lines 14-18) and the native theorem declarations with the mathematical macros (native preamble lines 154-217).\par\medskip
\noindent\textit{Allowed standard prerequisites (editorial).}\begin{itemize}
\item Big site (Sch/X)\_tau for a topology tau and full subcategories closed under coverings and products with the affine line. Same chapter, the big Zariski, etale, smooth, syntomic and fppf sites of a scheme and the hypotheses on the subcategories C\_1 and C\_2. Supplies the two sites between which the morphism of topoi g is formed and the hypotheses used in the proof.
\item Continuity and cocontinuity of the inclusion of the subcategories and the induced morphism of topoi. Same chapter, the morphism of topoi g : Sh(C\_1) -> Sh(C\_2) attached to the fully faithful continuous and cocontinuous inclusion. Supplies the morphism whose exceptional inverse image g\_! is computed.
\item The functor g\_! and its description as the sheafification of g\_\{p!\}. Modules on Sites, the definition of g\_! for a morphism of topoi and its presheaf-level description by a colimit. Supplies the presheaf whose sheafification is compared with O\_2 in the proof.
\end{itemize}\medskip
\ifdefined\hypertarget\hypertarget{sourcepointer-1}{}\fi
\noindent\textit{Source reference.} Original definition 04BF: sources/stacks-etale-cohomology/pinned-original-source.tex, source lines 7094--7094. Cross-chapter navigation pointer occurring inside the counted statement at source line 7094, where the author names the functor g\_! defined in another chapter. The target chapter is not staged in this card's source tree, so the pointer records the invoking line; it only resolves the reference for the reader, and no bibliography citation is involved.\par\medskip
\begin{lemma}
\label{lemma-compare-structure-sheaves}
Let $X$ be a scheme. Let
$\tau \in \{Zariski, \etale, smooth, syntomic, fppf\}$.
Let $\mathcal{C}_1 \subset \mathcal{C}_2 \subset (\Sch/X)_\tau$ be full
subcategories with the following properties:
\begin{enumerate}
\item For an object $U/X$ of $\mathcal{C}_t$,
\begin{enumerate}
\item if $\{U_i \to U\}$ is a covering of $(\Sch/X)_\tau$, then
$U_i/X$ is an object of $\mathcal{C}_t$,
\item $U \times \mathbf{A}^1/X$ is an object of $\mathcal{C}_t$.
\end{enumerate}
\item $X/X$ is an object of $\mathcal{C}_t$.
\end{enumerate}
We endow $\mathcal{C}_t$ with the structure of a site whose coverings are
exactly those coverings $\{U_i \to U\}$ of $(\Sch/X)_\tau$ with
$U \in \Ob(\mathcal{C}_t)$. Then
\begin{enumerate}
\item[(a)] The functor $\mathcal{C}_1 \to \mathcal{C}_2$
is fully faithful, continuous, and cocontinuous.
\end{enumerate}
Denote $g : \Sh(\mathcal{C}_1) \to \Sh(\mathcal{C}_2)$ the corresponding
morphism of topoi. Denote $\mathcal{O}_t$ the restriction of $\mathcal{O}$
to $\mathcal{C}_t$. Denote $g_!$ the functor of
Modules on Sites, Definition \ref{sites-modules-definition-g-shriek}.
\begin{enumerate}
\item[(b)] The canonical map $g_!\mathcal{O}_1 \to \mathcal{O}_2$
is an isomorphism.
\end{enumerate}
\end{lemma}

\begin{proof}
Assertion (a) is immediate from the definitions.
In this proof all schemes are schemes over $X$ and all morphisms of
schemes are morphisms of schemes over $X$. Note that $g^{-1}$ is
given by restriction, so that for an object $U$ of $\mathcal{C}_1$
we have $\mathcal{O}_1(U) = \mathcal{O}_2(U) = \mathcal{O}(U)$.
Recall that $g_!\mathcal{O}_1$ is the sheaf associated to the presheaf
$g_{p!}\mathcal{O}_1$ which associates to $V$ in $\mathcal{C}_2$ the group
$$
\colim_{V \to U} \mathcal{O}(U)
$$
where $U$ runs over the objects of $\mathcal{C}_1$ and the colimit is
taken in the category of abelian groups. Below we will use frequently
that if
$$
V \to U \to U'
$$
are morphisms with $U, U' \in \Ob(\mathcal{C}_1)$
and if $f' \in \mathcal{O}(U')$ restricts to $f \in \mathcal{O}(U)$,
then $(V \to U, f)$ and $(V \to U', f')$ define the same element of the
colimit. Also, $g_!\mathcal{O}_1 \to \mathcal{O}_2$ maps the element
$(V \to U, f)$ simply to the pullback of $f$ to $V$.

\medskip\noindent
Surjectivity. Let $V$ be a scheme and let $h \in \mathcal{O}(V)$.
Then we obtain a morphism $V \to X \times \mathbf{A}^1$ induced by $h$
and the structure morphism $V \to X$. Writing
$\mathbf{A}^1 = \Spec(\mathbf{Z}[x])$ we see the element
$x \in \mathcal{O}(X \times \mathbf{A}^1)$ pulls
back to $h$. Since $X \times \mathbf{A}^1$ is an object of $\mathcal{C}_1$
by assumptions (1)(b) and (2) we obtain the desired surjectivity.

\medskip\noindent
Injectivity. Let $V$ be a scheme. Let
$s = \sum_{i = 1, \ldots, n} (V \to U_i, f_i)$ be an element of the colimit
displayed above. For any $i$ we can use the morphism
$f_i : U_i \to X \times \mathbf{A}^1$
to see that $(V \to U_i, f_i)$ defines the same element of the colimit as
$(f_i : V \to X \times \mathbf{A}^1, x)$. Then we can consider
$$
f_1 \times \ldots \times f_n : V \to X \times \mathbf{A}^n
$$
and we see that $s$ is equivalent in the colimit to
$$
\sum\nolimits_{i = 1, \ldots, n}
(f_1 \times \ldots \times f_n : V \to X \times \mathbf{A}^n, x_i) =
(f_1 \times \ldots \times f_n : V \to X \times \mathbf{A}^n,
x_1 + \ldots + x_n)
$$
Now, if $x_1 + \ldots + x_n$ restricts to zero on $V$, then we see
that $f_1 \times \ldots \times f_n$ factors through
$X \times \mathbf{A}^{n - 1} = V(x_1 + \ldots + x_n)$. Hence we see
that $s$ is equivalent to zero in the colimit.
\end{proof}

\end{document}
